\section*{Capítulo \ref{ch:g6:water-underground} --- El agua bajo tierra}

\begin{solution}{exo:g6:water-underground:1}
\omterm{def:g6:water-underground:permeability}{Permeables}: \omterm{def:g3:stones-pebbles-sand:sand}{arena}, \omterm{def:g3:stones-pebbles-sand:pebble}{grava}, creta. \omterm{def:g6:water-underground:permeability}{Impermeables}: arcilla, granito sin
grietas, lutita.
\end{solution}

\begin{solution}{exo:g6:water-underground:2}
(a) \omterm{def:g6:water-underground:water-table}{nivel freático}; (b) \omterm{def:g6:water-underground:aquifer}{acuífero}; (c) \omterm{def:g6:water-underground:spring}{manantial}; (d) \omterm{def:g6:water-underground:spring}{pozo}.
\end{solution}

\begin{solution}{exo:g6:water-underground:3}
(a) Se evapora y vuelve al aire; corre por la superficie hasta los arroyos;
se infiltra en el suelo (\omterm{def:g6:water-underground:infiltration}{infiltración}). (b) En el aparcamiento: su
superficie dura casi no deja que el agua se infiltre, así que el agua corre
por encima. En el prado, el suelo y la hierba dejan que buena parte se
infiltre.
\end{solution}

\begin{solution}{exo:g6:water-underground:4}
(a) Falso: la arcilla es muy \omterm{def:g6:water-underground:porosity}{porosa} pero \omterm{def:g6:water-underground:permeability}{impermeable}, porque sus \omterm{def:g6:water-underground:porosity}{poros} son
diminutos. (b) Verdadero. (c) Falso: sigue la forma del terreno, de manera
más suave, y sube y baja con las estaciones. (d) Verdadero: alrededor del
30~\% de toda el \omterm{def:g2:water-on-earth:salt-fresh}{agua dulce} es \omterm{def:g6:water-underground:groundwater}{agua subterránea}, frente a alrededor de la
cuarta parte del uno por ciento en los \omterm{def:g1:landscapes-around-us:water}{lagos}.
\end{solution}

\begin{solution}{exo:g6:water-underground:5}
$10 \times 25 \div 100 = 2.5$: puedes echar 2.5~litros antes de que rebose.
$10 \times 22 \div 100 = 2.2$: saldrían 2.2~litros por el agujero.
\end{solution}

\begin{solution}{exo:g6:water-underground:6}
(a) $2 \times 1 \times 1 = 2$~metros cúbicos, es decir, 2000~litros.
(b) $2000 \times 20 \div 100 = 400$~litros. (c) Salen $2000 \times 18 \div 100 =
360$~litros; se quedan pegados $400 - 360 = 40$~litros.
\end{solution}

\begin{solution}{exo:g6:water-underground:7}
(a) El \omterm{def:g6:water-underground:spring}{manantial} está a unos 203~m, y el \omterm{def:g1:landscapes-around-us:water}{río}, a unos 132~m. (b) El
\omterm{def:g6:water-underground:spring}{manantial} está unos 70~m por encima del \omterm{def:g1:landscapes-around-us:water}{río}: su agua baja por la pendiente
hasta el \omterm{def:g1:landscapes-around-us:water}{río}. (c) Dentro de la ladera de creta, el \omterm{def:g6:water-underground:water-table}{nivel freático} está por
debajo del suelo; solo se encuentra con la superficie donde la base de la
creta, apoyada sobre la lutita \omterm{def:g6:water-underground:permeability}{impermeable}, sale a la pendiente.
\end{solution}

\begin{solution}{exo:g6:water-underground:8}
(a) Después de la estación lluviosa: $401 - 336 = 65$~m; después de la
estación seca: $401 - 314 = 87$~m. (b) $336 - 314 = 22$~m.
\end{solution}

\begin{solution}{exo:g6:water-underground:9}
Más profundo. El \omterm{def:g6:water-underground:water-table}{nivel freático} sigue la forma del terreno, pero de manera
más suave: en lo alto de la \omterm{def:g1:landscapes-around-us:landform}{colina} está más alto que bajo el \omterm{def:g6:water-underground:spring}{pozo}~A, pero
el suelo sube mucho más, así que el agua está más lejos de la superficie
(en el corte, a más de 110~m de profundidad).
\end{solution}

\begin{solution}{exo:g6:water-underground:10}
Los \omterm{def:g6:water-underground:porosity}{poros} de la arcilla son diminutos: el agua se queda pegada a las
láminas y no puede atravesar la arcilla hacia el \omterm{def:g6:water-underground:spring}{pozo}. La arcilla es \omterm{def:g6:water-underground:porosity}{porosa}
pero \omterm{def:g6:water-underground:permeability}{impermeable}, así que solo sale un poco de agua (20~litros de 500).
\end{solution}

\begin{solution}{exo:g6:water-underground:11}
(a) $30.1 \div 0.26 \approx 116$: más de cien veces más. (b)
$100 - 69 = 31$: queda alrededor del 31~\%, y solo el \omterm{def:g6:water-underground:groundwater}{agua subterránea}
supone unos 30 de ellos. (c) Es, con mucho, la mayor reserva de \omterm{def:g2:water-on-earth:salt-fresh}{agua dulce}
líquida; el hielo está congelado y casi siempre lejos, en las regiones
polares y en las \omterm{def:g1:landscapes-around-us:landform}{montañas}. Y hay \omterm{def:g6:water-underground:groundwater}{agua subterránea} casi en todas partes bajo
nuestros pies.
\end{solution}

\begin{solution}{exo:g6:water-underground:12}
Los dos \omterm{def:g6:water-underground:spring}{pozos} sacan agua del mismo \omterm{def:g6:water-underground:aquifer}{acuífero} de creta. La bomba potente sacó
el agua más deprisa de lo que la \omterm{def:g4:clouds-rain-snow:precipitation}{lluvia} vuelve a llenar la creta, así que
el \omterm{def:g6:water-underground:water-table}{nivel freático} bajó en toda la zona, por debajo del fondo del segundo
\omterm{def:g6:water-underground:spring}{pozo}, que se secó. Los granjeros podrían ponerse de acuerdo para bombear
menos, sobre todo en verano, compartir un \omterm{def:g6:water-underground:spring}{pozo} o hacer más hondo el segundo
\omterm{def:g6:water-underground:spring}{pozo}; sobre todo, no deberían sacar más agua de la que devuelve la \omterm{def:g4:clouds-rain-snow:precipitation}{lluvia}.
\end{solution}

\begin{solution}{pb:g6:water-underground:1}
\textbf{1.} La granja está sobre creta; debajo hay lutita. La creta es
\omterm{def:g6:water-underground:permeability}{permeable}; la lutita, no.

\textbf{2.} La \omterm{def:g4:clouds-rain-snow:precipitation}{lluvia} se infiltra en la creta y se detiene sobre la
lutita. Donde la base de la creta sale a la pendiente, el \omterm{def:g6:water-underground:water-table}{nivel freático} se
encuentra con el suelo y el agua sale: es el \omterm{def:g6:water-underground:spring}{manantial}.

\textbf{3.} $401 - 336 = 65$~m por debajo del suelo.

\textbf{4.} El fondo del \omterm{def:g6:water-underground:spring}{pozo} está a $401 - 100 = 301$~m, por debajo del
\omterm{def:g6:water-underground:water-table}{nivel freático}, que está a 314~m: el \omterm{def:g6:water-underground:spring}{pozo} sigue llegando al agua, por $314 - 301 =
13$~m.

\textbf{5.} $100 \times 50 \times 2 = 10\,000$~metros cúbicos.

\textbf{6.} $10\,000 \times 20 \div 100 = 2000$~metros cúbicos.

\textbf{7.} Salen $10\,000 \times 18 \div 100 = 1800$~metros cúbicos; se
quedan pegados $2000 - 1800 = 200$~metros cúbicos.

\textbf{8.} $1800 \times 1000 = 1\,800\,000$~litros.

\textbf{9.} $18 \times 1000 = 18\,000$~litros al día.

\textbf{10.} $18 \times 7 = 126$~metros cúbicos a la semana.

\textbf{11.} La \omterm{def:g4:clouds-rain-snow:precipitation}{lluvia} de la estación húmeda se infiltra en la creta y la
vuelve a llenar cada año (mientras las granjas no bombeen más de lo que
devuelve la \omterm{def:g4:clouds-rain-snow:precipitation}{lluvia}).

\textbf{12.} $1800 \div 18 = $~\textbf{100 días}: algo más de tres meses
sin ninguna \omterm{def:g4:clouds-rain-snow:precipitation}{lluvia}.
\end{solution}
